\section{Motivation}
\label{sec:motivation}

High-Performance Computing (HPC) has become an essential infrastructure in modern scientific research, engineering simulations, and large-scale data processing. HPC clusters enable researchers to tackle computationally intensive problems that would be infeasible on conventional computing systems. From climate modeling and genomics to machine learning and financial simulations, HPC systems provide the computational power necessary to drive innovation across multiple domains.

Slurm Workload Manager, originally introduced as the Simple Linux Utility for Resource Management, has emerged as a widely deployed workload manager for HPC clusters \cite{yoo2003slurm, slurm2024docs}. SchedMD reports that Slurm is used at government laboratories, universities, and companies worldwide, and that it manages workloads for over half of the top 10 systems in the TOP500 list \cite{schedmd2024slurm}. Slurm provides sophisticated mechanisms for job scheduling, resource allocation, and cluster monitoring, enabling efficient utilization of computational resources across thousands of nodes. The system's flexibility and scalability have made it the preferred choice for many academic institutions and commercial organizations operating large-scale computing facilities.

Despite its power and flexibility, HPC user-support literature highlights recurring needs for training, documentation, and assistance when researchers interact with complex computing environments \cite{sabin2020hpc}. In Slurm-based clusters, this challenge is reflected in a broad command-line interface (CLI) comprising numerous commands, each with extensive sets of options and parameters \cite{slurm2024docs}. Users must master commands such as \texttt{squeue}, \texttt{sinfo}, \texttt{sbatch}, \texttt{scontrol}, and \texttt{sacct}, along with their associated flags and output formats. This complexity is compounded by the need to write job submission scripts in Bash, which requires additional technical knowledge beyond the domain expertise of many researchers.

Recent advances in Large Language Models (LLMs) have demonstrated remarkable capabilities in natural language understanding and generation \cite{brown2020language, openai2023gpt4}. These models can interpret complex user queries, reason about multi-step tasks, and generate appropriate responses or actions. The emergence of LLM-based agents---systems that combine language models with tool-calling capabilities---has opened new possibilities for creating intelligent interfaces to complex systems \cite{xi2023rise, schick2023toolformer}.

The Model Context Protocol (MCP), developed by Anthropic, provides a standardized approach for connecting LLMs to external tools and data sources \cite{anthropic2024mcp, mcpspec2024}. MCP enables the creation of modular, reusable integrations that can expose system functionality to AI agents through explicit tool interfaces. This architectural pattern separates the concerns of tool implementation from agent logic, facilitating the development of sophisticated AI-assisted applications.

However, deploying commercial LLM APIs in HPC environments raises fundamental concerns. Cluster workloads expose sensitive metadata---user submission scripts, resource allocation patterns, job failure diagnostics, and access-control policies---that institutional data governance policies prohibit from traversing external networks. Furthermore, recurring per-query API costs scale linearly with cluster user-base size. These constraints motivate the deployment of locally-hosted open-weight models that preserve data sovereignty and eliminate ongoing costs, but at the expense of model capacity: a 14-billion parameter model cannot match the tool-selection accuracy of commercial frontier models when presented with a large number of available tools simultaneously.

This research is motivated by the opportunity to bridge the gap between Slurm's powerful capabilities and user accessibility through an intelligent conversational interface that operates entirely on local infrastructure. By leveraging LLM agents connected to Slurm through MCP, and by partitioning the tool surface across specialized sub-agents to compensate for the capacity constraints of local models, we create a system that allows users to manage cluster resources using natural language queries while maintaining full data sovereignty.


